%%%%%%%%%%%%%%%%%%%%%%%%%%%%%%%%%%%%%%%%%%%%%%%%%%%%%%%%%%%%%%%%%%%%%%%%%%%%%%
%  Unit IV appendices
%%%%%%%%%%%%%%%%%%%%%%%%%%%%%%%%%%%%%%%%%%%%%%%%%%%%%%%%%%%%%%%%%%%%%%%%%%%%%%

\startappendices

\chapter{Syllabus-to-Chapter Concordance}

This appendix exists so that coverage can be \emph{audited} rather than assumed.
The left-hand column reproduces the Unit~IV syllabus descriptor phrase by phrase,
in the order the official document prints it. The right-hand column gives the
chapter and section that carries it. Every phrase resolves to a location, and no
numbered chapter of this volume contains material that is not traceable to a
phrase in this table.

\begin{mmlongtable}{Concordance for Unit IV}{tab:concordance4}%
{@{}P{0.36\textwidth}P{0.57\textwidth}@{}}
\rowcolor{ocre}\thd{Syllabus phrase} & \thd{Where it is covered}\\
\midrule
\endfirsthead
\rowcolor{ocre}\thd{Syllabus phrase} & \thd{Where it is covered}\\
\midrule
\endhead
\multicolumn{2}{@{}l}{\bfseries\color{ocredark} First half --- Web Design}\\
\midrule
\rlab{Web design, optimization of websites} & Chapter~1 entire. What web design is
\S\,1.1; UI and UX \S\,1.2; website structure \S\,1.3; website optimization
defined \S\,1.4; the eight techniques \S\,1.5; image optimization \S\,1.6;
testing and optimisation \S\,1.7.\\
\rowcolor{tablealt}
\rlab{Publishing a basic website} & Chapter~2 entire. The nine official steps
\S\,2.1; each step expanded, including the pre-launch technical setup \S\,2.2.\\
\rlab{User-centered Design and Website Optimization} & Chapter~3 entire. What UCD
is \S\,3.1; the seven principles \S\,3.2; the UCD process \S\,3.3; what visitors
look for \S\,3.4. Website optimization itself is carried in \S\,1.4--1.7.\\
\rowcolor{tablealt}
\rlab{Design Principles and Website Copy} & Chapter~4 entire. The core design
principles \S\,4.1; Krug's first law and the usability attributes \S\,4.2; how
people use the web \S\,4.3; billboard design \S\,4.4; copy principles \S\,4.5;
design and copy as one system \S\,4.6.\\
\rlab{Website Metrics \& Developing Insight} & Chapter~5 entire. Key metrics
\S\,5.1; how metrics aid optimisation \S\,5.2; from data to decision \S\,5.3;
user behaviour analysis \S\,5.4.\\
\midrule
\multicolumn{2}{@{}l}{\bfseries\color{ocredark} Second half --- Mobile Marketing}\\
\midrule
\rowcolor{tablealt}
\rlab{Difference between mobile advertising and marketing} & Chapter~6 entire.
What mobile marketing is \S\,6.1; advertising against marketing \S\,6.2; mobile
web against app \S\,6.3; benefits \S\,6.4; limitations \S\,6.5; the seven
campaign characteristics \S\,6.6; geo-targeting \S\,6.7.\\
\rlab{Utilizing mobile marketing for sales promotions} & Chapter~7 entire. Push
notifications \S\,7.1; ten promotion techniques \S\,7.2; advertising formats
\S\,7.3.\\
\rowcolor{tablealt}
\rlab{Online applications} & Chapter~8 entire. Why apps matter \S\,8.1; business
models \S\,8.2; the three creation rules \S\,8.3; the feature checklist \S\,8.4;
creating a mobile website \S\,8.5; mobile analytics \S\,8.6.\\
\end{mmlongtable}

\begin{keybox}[Two Notes on Coverage]
\textbf{Deliberate overlap with Volume~1.} Krug's first law, the seven usability\ix{usability}
attributes and the landing-page material appear in Unit~I under ``Optimize
Website UX \& Landing Pages'' and again here under ``Design Principles''. Both
descriptors name them, so both volumes carry them. Chapter~4 gives the design-side
treatment; Volume~1 Chapter~9 gives the content-marketing side.

\smallskip
\textbf{No off-syllabus appendix.} As in Volume~3, every topic in the source
material for Unit~IV maps to a descriptor phrase, so nothing had to be set aside.
\end{keybox}

\chapter{Previous-Year Question Bank --- Unit IV}

Every question in this appendix is reproduced from an actual RVCE Marketing
Management Continuous Internal Evaluation or Semester End Examination paper for
this course. Answers and pointers follow the official schemes of valuation where
those were released.

\section*{Part A --- Two-Mark Questions}
\addcontentsline{toc}{section}{Part A --- Two-Mark Questions}

\begin{enumerate}[leftmargin=2.2em,itemsep=1.1ex]

\item \textbf{What is Website Optimization?} \pyqr{CIE-III, Jun 2026 --- 2 M}\\
      The process of improving website performance, speed, usability and search
      visibility in order to increase traffic and conversions.
      \hfill\emph{See} \S\,1.4.

\item \textbf{Mention any two principles of website design.}
      \pyqr{CIE-III, Jun 2026 --- 2 M}\\
      Simplicity and consistency --- these principles improve readability and user
      experience. \hfill\emph{See} \S\,4.1.

\item \textbf{Differentiate Mobile Advertising and Mobile Marketing. / What is the
      key difference between mobile advertising and mobile marketing? / Define
      mobile advertising in one sentence. How does it differ from mobile
      marketing?}
      \pyqr{CIE-III, Jun 2026; Improvement CIE, 4 Jun 2025; Improvement CIE, 28 Aug 2024 --- 2 M}\\
      Mobile advertising focuses mainly on paid promotional messages delivered to
      mobile devices. Mobile marketing includes overall customer engagement
      through SMS, apps, social media and mobile websites. In one line: mobile
      advertising is a paid subset of the broader mobile marketing strategy.
      \hfill\emph{See} \S\,6.2.

\item \textbf{Mention two key principles of User-Centered Design in website
      development.} \pyqr{Improvement CIE, 4 Jun 2025 --- 2 M}\\
      (1) Design based on an explicit understanding of the users, their tasks and
      their environment. (2) Iterative design that is continuously refined
      through user involvement and user-centred evaluation.
      \hfill\emph{See} \S\,3.1.

\item \textbf{How can website metrics help in optimizing a website for marketing
      effectiveness?} \pyqr{Improvement CIE, 4 Jun 2025 --- 2 M}\\
      They locate problems precisely, rank fixes by impact, prove or disprove
      changes through testing, reallocate budget to converting channels, and
      expose segment-specific failures hidden by averages.
      \hfill\emph{See} \S\,5.2.

\item \textbf{State two key elements that visitors commonly expect to see on a
      well-optimized landing page.} \pyqr{Improvement CIE, 4 Jun 2025 --- 2 M}\\
      A clear benefit-led headline matching the advertisement or link that brought
      them there, and a single prominent call-to-action button. Also acceptable:
      trust signals such as testimonials or ratings. \hfill\emph{See} \S\,3.4.

\item \textbf{How does image optimization contribute to overall website
      performance?} \pyqr{Improvement CIE, 28 Aug 2024 --- 2 M}\\
      It reduces page weight and load time, improving Core Web Vitals\ix{Core Web Vitals} and
      lowering bounce\ix{bounce} rate; improves the mobile experience on slow connections;
      improves SEO through descriptive file names and ALT text; and improves
      accessibility. \hfill\emph{See} \S\,1.6.

\item \textbf{List any four key metrics used to measure the success of a website.}
      \pyqr{Improvement CIE, 28 Aug 2024 --- 2 M}\\
      Traffic (users or sessions), bounce rate, average session duration and
      conversion rate. Also acceptable: pages per session, page load time,
      traffic sources, new against returning visitors.
      \hfill\emph{See} \S\,5.1.

\item \textbf{How can push notifications be used in mobile marketing to enhance
      sales promotions?} \pyqr{Improvement CIE, 28 Aug 2024 --- 2 M}\\
      They deliver time-critical offers instantly to the installed base, create
      urgency for flash sales, recover abandoned carts, trigger location-based
      coupons when a user nears a store, and re-engage lapsed users --- all
      through an owned channel that no algorithm throttles.
      \hfill\emph{See} \S\,7.1.

\item \textbf{How can user-friendly design impact the overall effectiveness of a
      website?} \pyqr{Improvement CIE, 28 Aug 2024 --- 2 M}\\
      It reduces cognitive effort so that visitors can complete their task, which
      lowers bounce rate, increases time on site and pages per session, raises
      conversion rate, builds trust in the brand, encourages return visits, and
      --- because search engines now mimic human preferences --- indirectly
      improves organic ranking. \hfill\emph{See} \S\,4.2.

\item \textbf{Why is it important to integrate SEO in website development?}
      \pyqr{SEE, Jun/Jul 2025 --- 2 M}\\
      Architecture, URLs, rendering method, speed, heading structure and analytics
      are all fixed at build time. Retro-fitting SEO is expensive, forces URL
      changes that risk losing ranking authority, and often leaves staging
      \texttt{noindex} directives live in production.
      \hfill\emph{See} \S\,2.2, and Volume~2 \S\,8.2.

\item \textbf{What is the role of geo-targeting in mobile marketing?}
      \pyqr{SEE, Jun/Jul 2025 --- 2 M}\\
      Geo-targeting delivers content and offers based on the user's location,
      making messages locally relevant and timely --- driving footfall through
      geo-fenced offers, enabling ``near me'' discovery and click-to-map
      directions, supporting location-specific pricing and inventory, and
      improving ROI by eliminating spend outside the serviceable area.
      \hfill\emph{See} \S\,6.7.

\item \textbf{List any two aspects of branding and consistency in web design.}
      \pyqr{SEE, Oct/Nov 2023 --- 2 M}\\
      (1) A consistent visual identity --- the same logo, colour palette,
      typography, imagery style and button design applied on every page. (2) A
      consistent tone of voice and terminology in all copy, including microcopy,
      error messages and CTA labels, so that the brand reads as one coherent
      personality. \hfill\emph{See} \S\,4.1.

\item \textbf{List any two differences between Mobile Web marketing and Mobile App
      marketing.} \pyqr{SEE, Oct/Nov 2023 --- 2 M}\\
      (1) Mobile web is accessed through a browser with no installation and
      therefore has broader reach, whereas an app requires download and
      installation and reaches a narrower but far more engaged audience. (2) Apps
      have full push notification and device-feature access --- camera, GPS,
      biometrics --- and can work offline, whereas mobile web has limited push
      capability and device access. \hfill\emph{See} \S\,6.3.

\item \textbf{What is the impact of user behavior analysis on website
      optimization?} \pyqr{SEE, Oct/Nov 2023 --- 2 M}\\
      It shows what users actually do rather than what is assumed --- identifying
      friction points, ignored content, invisible calls to action, rage clicks and
      true scroll depth\ix{scroll depth} --- and generates the hypotheses that A/B testing then
      validates. Without it, optimisation is guesswork.
      \hfill\emph{See} \S\,5.4.

\end{enumerate}

\section*{Part B --- Eight- and Ten-Mark Questions}
\addcontentsline{toc}{section}{Part B --- Eight- and Ten-Mark Questions}

\begin{enumerate}[leftmargin=2.2em,itemsep=1.4ex]

\item \textbf{Explain the concept of Web Design and the importance of website
      optimization.} \pyqr{CIE-III, Jun 2026 --- 10 M}\\
      Define web design covering layout, colour, typography, graphics, navigation
      and responsiveness. Then give the six importance points, the four types of
      optimization --- technical, content, SEO, mobile --- and the three benefits.
      \hfill\emph{See} \S\,1.1 and \S\,1.4.

\item \textbf{A start-up company notices that visitors leave their website within
      a few seconds. As a web designer, explain how you would optimize the website
      to improve user engagement and conversion rate.}
      \pyqr{CIE-III, Jun 2026 --- 10 M}\\
      Give the eight numbered techniques --- speed, responsive design\ix{responsive design}, navigation,
      content quality, strong calls to action, SEO, bounce-rate reduction and
      metric monitoring --- each with its concrete sub-actions. Strengthen the
      answer by opening with a diagnosis step (heatmaps and session recordings to
      find \emph{why} they leave) and closing with A/B testing\ix{A/B testing} to validate the
      fixes. \hfill\emph{See} \S\,1.5.

\item \textbf{A college plans to publish its departmental website for admissions
      and student activities. Explain the complete process involved in designing
      and publishing the website.} \pyqr{CIE-III, Jun 2026 --- 10 M}\\
      Give the nine official steps --- requirement analysis, website design,
      content development, domain registration, web hosting, website development,
      testing, publishing, maintenance --- expanded with the practitioner detail,
      including the pre-launch technical setup. \hfill\emph{See} Chapter~2.

\item \textbf{Discuss what website visitors typically look for when they land on a
      business webpage. Suggest effective design and content strategies to meet
      those expectations.} \pyqr{Improvement CIE, 4 Jun 2025 --- 10 M}\\
      Use the ten visitor questions with the paired delivery strategy for each,
      then add the design principles and the copy principles as the systematic
      response. \hfill\emph{See} \S\,3.4, \S\,4.1 and \S\,4.5.

\item \textbf{Design a user-centered website layout for a new e-commerce startup
      targeting young urban professionals. How would you ensure optimal user
      experience and engagement? Include key design principles, optimization
      strategies, and website copy suggestions.}
      \pyqr{Improvement CIE, 4 Jun 2025 --- 10 M}\\
      \emph{Research first:} persona --- 25 to 35, urban, mobile-first, time-poor,
      price-aware but convenience-driven; primary task --- find and buy quickly on
      a phone during a commute.
      \emph{Layout:} sticky header with search and cart; a hero section with one
      benefit headline and one call to action\ix{call to action}; a three-column category grid; a
      social-proof band with ratings; a personalised ``recommended for you'' row;
      a simple footer with trust and policy links.
      \emph{Key design principles applied:} mobile-first responsive layout; visual
      hierarchy placing the search bar and primary CTA highest; simplicity with
      generous whitespace; consistency in button style and colour; one accent
      colour reserved for actions; 16-pixel or larger type; tap targets of at
      least 44 pixels; accessibility through ALT text and contrast.
      \emph{Optimization strategies:} compressed modern-format images and lazy
      loading for a sub-three-second load; one-page guest checkout with a
      four-field form and saved payment; a progress indicator in checkout; an
      exit-intent offer; heatmaps and session recordings to locate friction; A/B
      tests on the hero headline and CTA.
      \emph{Copy suggestions:} headline ``Everything for your week, delivered by
      tomorrow''; sub-line ``Free delivery over \rupee 499. Returns in one tap.'';
      CTA ``Shop today's deals'' rather than ``Submit''; microcopy under the form
      reading ``Takes 30 seconds --- no account needed''; benefit-led product copy
      with specific claims.
      \emph{Measurement:} conversion rate, checkout completion, mobile bounce rate
      and average order value. \hfill\emph{See} Chapters~3, 4 and~5.

\item \textbf{A travel agency has created a landing page for a limited-time offer
      on holiday packages. Despite a high number of clicks on the ads leading to
      the landing page, the conversion rate is disappointingly low. Evaluate the
      possible reasons for the low conversion rate. Suggest optimizations for key
      elements such as headlines, call-to-action buttons, and visuals.}
      \pyqr{Improvement CIE, 28 Aug 2024 --- 10 M}\\
      \emph{Possible reasons:} poor message match\ix{message match} --- the advertisement promised
      something the page does not repeat; a vague headline that does not restate
      the offer; the offer and price are below the fold; too many competing calls
      to action, or none that is visually dominant; a long enquiry form demanding
      unnecessary details; no urgency despite the offer being time-limited;
      missing trust signals --- no reviews, no certifications, no cancellation
      policy; slow page load, especially on mobile; generic stock imagery that does
      not show the actual destinations; hidden total cost or unclear inclusions;
      and navigation links that let visitors wander away from the page.

      \emph{Optimisations --- headlines:} restate the exact advertisement promise
      with specifics, such as ``Kerala backwaters, 4 nights from \rupee 18,999 ---
      book by Sunday'', and add a sub-headline listing what is included.
      \emph{CTA buttons:} one dominant, high-contrast accent-coloured button
      repeated above and below the fold, with active benefit copy such as
      ``Reserve my package'' instead of ``Submit'', plus reassurance microcopy
      (``Free cancellation up to 7 days'').
      \emph{Visuals:} replace stock photography with real destination and customer
      photographs, add a short itinerary video, and show a countdown timer for the
      deadline.
      \emph{Additional:} cut the form to name, phone and travel date; add
      testimonials and ratings adjacent to the CTA; remove the site navigation from
      the landing page; compress images for a sub-three-second mobile load; and
      A/B test the headline and CTA.
      \emph{Measurement:} conversion rate, form-start to form-completion ratio,
      scroll depth, and session recordings to verify the fix.
      \hfill\emph{See} \S\,3.4, \S\,4.5 and \S\,5.3.

\item \textbf{How do online applications serve as effective tools for mobile
      marketing? What are the key features that businesses should incorporate into
      their apps to ensure they are maximizing marketing potential?}
      \pyqr{Improvement CIE, 28 Aug 2024 --- 10 M}\\
      The eight reasons apps are effective, the four business models, the three
      app-creation rules, and the eleven-feature checklist.
      \hfill\emph{See} Chapter~8.

\item \textbf{A luxury cosmetics brand is planning to launch a flash sale
      exclusively through mobile devices. How would you design a mobile marketing
      campaign that not only drives awareness of the sale but also maximizes
      conversion rates during the limited time?}
      \pyqr{Improvement CIE, 28 Aug 2024 --- 10 M}\\
      \emph{Pre-launch (awareness, seven days out):} teaser push notifications and
      SMS to the segmented installed base; Reels and short-video influencer
      previews; an in-app countdown banner; a ``notify me'' opt-in that captures
      intent and grants push permission; paid social and mobile display
      retargeting to past purchasers and lookalikes.
      \emph{Launch (conversion window):} a timed push notification at each
      segment's peak activity hour, deep-linked straight to the sale screen;
      app-exclusive early access one hour before the website, to drive installs; a
      visible countdown timer on the sale screen; personalised product rows based
      on browsing and purchase history; limited-stock indicators to create
      scarcity; one-tap checkout with saved cards and wallet support.
      \emph{During (maximising conversion):} an abandoned-cart push within 20
      minutes; geo-fenced offers near flagship stores with click-to-map; live chat
      support in-app; free-shipping threshold prompts to raise average order value.
      \emph{Post-launch:} thank-you and review-request push; a win-back offer to
      browsers who did not buy; and a wallet pass for the next drop.
      \emph{Technical requirements:} sub-three-second load, server capacity for the
      traffic spike, and a fully responsive mobile-first design.
      \emph{Measurement:} push open rate, sale-screen sessions, add-to-cart rate,
      checkout completion, conversion rate, average order value, revenue per push,
      opt-out rate, and app installs during the window.
      \hfill\emph{See} Chapters~7 and~8.

\item \textbf{Assume you are launching a promotional campaign for a new product
      using mobile marketing. Prepare a detailed plan including mobile advertising
      channels, push notifications, app-based promotions, and integration with
      social media.} \pyqr{Improvement CIE, 4 Jun 2025 --- 10 M}\\
      \emph{Objectives and audience:} define a SMART\ix{SMART} objective and the mobile
      persona, including device, active hours and preferred platforms.
      \emph{Mobile advertising channels:} mobile search advertisements on
      high-intent queries with location extensions; in-app banner and native
      advertisements on relevant apps; rewarded video for reach; interstitials at
      natural app breakpoints; click-to-call advertisements for enquiry-driven
      products; click-to-map advertisements to drive footfall; and expandable
      rich-media advertisements for the brand story.
      \emph{Push notifications:} a sequenced plan --- teaser, launch, reminder,
      last chance --- each segmented, deep-linked, frequency-capped and A/B tested,
      with permission requested only after the user has experienced value.
      \emph{App-based promotions:} app-exclusive early access and pricing, in-app
      coupons, gamified spin-to-win, loyalty points, referral codes, wallet passes,
      and one-tap checkout.
      \emph{Social media integration:} Reels and Stories with the campaign
      hashtag; influencer seeding; shoppable posts and product tags; a messaging
      broadcast for order updates and offers; user-generated content\ix{user-generated content} contests; and
      retargeting audiences built from app and site events.
      \emph{Cross-channel integration:} QR codes on packaging and print linking to
      the mobile landing page; consistent creative and offer across all
      touchpoints; ``think mobile first'', with mobile as the connective tissue
      between all media.
      \emph{Measurement:} installs and cost per install, push open rate, in-app
      conversion rate, coupon redemption, return on advertising spend, and lifetime
      value by acquisition source. \hfill\emph{See} Chapters~6, 7 and~8.

\item \textbf{Discuss the relationship between design principles and website copy.
      How can visual design elements and well-crafted content work together to
      enhance user experience and achieve website objectives? Include examples of
      effective design and copy integration.}
      \pyqr{Improvement CIE, 28 Aug 2024 --- 10 M}\\
      The principle-by-principle table, followed by the worked examples of
      successful integration --- hero section, pricing table, testimonial block ---
      and of failure: a beautiful layout with vague copy, or strong copy set in
      unreadable type. \hfill\emph{See} \S\,4.6.

\item \textbf{What is the role of user Interface and user Experience designs in web
      design for marketing? How can businesses ensure a positive user Journey on
      their websites?} \pyqr{SEE, Oct/Nov 2023, Q7a --- 8 M}\\
      The UI and UX comparison table and the eight ways to ensure a positive user
      journey. \hfill\emph{See} \S\,1.2.

\item \textbf{Describe the importance of testing and optimization in web design for
      Marketing. What tools and methods can be used to analyze and improve Website
      performance?} \pyqr{SEE, Oct/Nov 2023, Q7b --- 8 M}\\
      The case for testing and the method-and-tool table.
      \hfill\emph{See} \S\,1.7.

\item \textbf{Discuss briefly on Push Notification in Mobile Marketing. Explain
      how can businesses use them to engage and retain Mobile App users? Provide
      examples of effective push notification strategies.}
      \pyqr{SEE, Oct/Nov 2023, Q8a --- 8 M}\\
      The definition, the seven engagement and retention uses, and the eight
      effective strategies --- with examples such as a 20-minute abandoned-cart
      reminder, a geo-fenced coupon on store approach, and a 30-day win-back
      message. \hfill\emph{See} \S\,7.1.

\item \textbf{Explain the significance of mobile analytics in mobile marketing
      campaigns. What key metrics should marketers track to measure the success of
      their mobile strategies?} \pyqr{SEE, Oct/Nov 2023, Q8b --- 8 M}\\
      The significance argument and the seven metric groups, from acquisition
      through to technical health. \hfill\emph{See} \S\,8.6.

\item \textbf{Evaluate the effectiveness of mobile marketing in facing market
      changes due to rapid technological advancements.}
      \pyqr{SEE, Jun/Jul 2025, Q8a --- 8 M}\\
      \emph{Where mobile marketing is effective:} it is inherently adaptive
      because it is digital and measurable --- campaigns can be tested and
      adjusted continuously; it reaches consumers wherever they are, at any hour;
      it spans the entire funnel from discovery to purchase; it costs less than
      traditional advertising; it produces first-party behavioural data that
      survives the loss of third-party cookies; and it enables geo-targeting,
      personalisation and instant response that no traditional channel can match.

      \emph{Where it is strained by technological change:} platform fragmentation
      across devices, screen sizes and operating systems makes a single campaign
      difficult; privacy shifts --- app tracking transparency, cookie deprecation,
      GDPR and do-not-disturb regulation --- have degraded attribution and
      targeting; ad fatigue and rising uninstall rates limit push frequency; and
      the pace of format change forces continual re-skilling.

      \emph{Evaluation:} mobile marketing remains the most resilient channel
      precisely because its measurability lets it adapt faster than the disruption
      --- but only for organisations that build an owned first-party asset, namely
      an app and a permission-based list, rather than renting reach from
      platforms. \hfill\emph{See} \S\,6.4, \S\,6.5 and \S\,6.6.

\item \textbf{Create a strategy to use new digital channels like mobile apps,
      chatbots, and social media integration to strengthen a brand's global online
      presence.} \pyqr{SEE, Jun/Jul 2025, Q8b --- 8 M}\\
      \emph{Mobile app as the owned hub:} a single global app with localised
      content, currency and language; personalised home screens; push notification
      consent captured after first value; loyalty and wallet integration; offline
      access; and full event instrumentation feeding a global analytics view.

      \emph{Chatbots as the always-on layer:} deployed on the website, in the app
      and on messaging platforms; handling FAQs, order tracking, returns and
      product discovery in multiple languages across all time zones; escalating to
      human agents with full context; and capturing conversational data that feeds
      product and content decisions.

      \emph{Social media integration:} a hub-and-spoke model with a global brand
      account and regional handles; platform-native content rather than
      cross-posting; shoppable posts and product tags linking back to the app;
      regional influencer partnerships for cultural credibility; and social
      listening for share of voice\ix{share of voice} and sentiment by market.

      \emph{Cross-channel integration:} one customer profile unified across app,
      chat, social, e-mail and web, so that the conversation continues seamlessly
      --- the omnichannel synchronisation principle of Volume~2 \S\,1.3.

      \emph{Governance:} a central brand system with regional flexibility;
      consistent visual identity and tone; and compliance with local privacy and
      messaging regulation.

      \emph{Measurement:} a global dashboard tracking reach, engagement, app
      installs and retention, chatbot resolution rate, cross-channel conversion and
      lifetime value by market. \hfill\emph{See} Chapters~6 and~8.

\end{enumerate}

\section*{How Unit IV Is Examined}
\addcontentsline{toc}{section}{How Unit IV Is Examined}

\begin{keybox}[The Pattern Across Papers]
In the Semester End Examination, Unit~IV is Questions~7 and~8 with internal
choice --- 16 marks in two eight-mark sub-divisions --- plus two or three Part~A
objective questions. The third Continuous Internal Evaluation test and the
Improvement CIE both draw heavily on it.

\smallskip
The highest-probability short question by a wide margin is the \textbf{mobile
advertising against mobile marketing} distinction, which has appeared in three
separate papers. The highest-probability ten-mark topics are the
\textbf{website optimisation scenario} (visitors leaving within seconds), the
\textbf{nine-step publishing process}, \textbf{what visitors look for} on a
business page, the \textbf{design-and-copy relationship}, \textbf{push
notifications}, and \textbf{online applications}.

\smallskip
Note that Unit~IV questions are heavily scenario-based --- a start-up, a college
department, a travel agency, a cosmetics brand, an e-commerce startup. Answer them
by naming the framework and then applying it to the specific business, not by
discussing web design in general.
\end{keybox}

\chapter{Glossary of Key Terms}

\begin{mmlongtable}{Key terms for rapid revision}{tab:glossary4}%
{@{}P{0.26\textwidth}P{0.67\textwidth}@{}}
\rowcolor{ocre}\thd{Term} & \thd{One-line meaning}\\
\midrule
\endfirsthead
\rowcolor{ocre}\thd{Term} & \thd{One-line meaning}\\
\midrule
\endhead
\rlab{Above the fold} & The area of a page visible without scrolling\\
\rowcolor{tablealt}
\rlab{ASO} & App store optimisation --- ranking an app in app store search\\
\rlab{Click-to-Call} & A mobile advertisement format with a tappable phone number;
lifts click-through by 6--8\%\\
\rowcolor{tablealt}
\rlab{Click-to-Map} & A geo-targeted format that opens a map to the nearest
store\\
\rlab{Core Web Vitals} & The page-experience metrics LCP, INP and CLS\\
\rowcolor{tablealt}
\rlab{CPI / LTV} & Cost per install / customer lifetime value\\
\rlab{DAU / MAU} & Daily and monthly active users; their ratio is
``stickiness''\\
\rowcolor{tablealt}
\rlab{Domain registration} & Buying the site's address\\
\rlab{Geo-fencing} & Triggering an action when a user enters a defined geographic
boundary\\
\rowcolor{tablealt}
\rlab{Geo-targeting} & Delivering content or offers based on the user's geographic
location\\
\rlab{Happy talk} & Empty welcome text that says nothing; should be deleted\\
\rowcolor{tablealt}
\rlab{Interstitial} & A full-screen advertisement displayed within an app\\
\rlab{Krug's first law} & ``Don't make me think'' --- the interface must be
self-evident\\
\rowcolor{tablealt}
\rlab{M-commerce} & Commercial transactions conducted through a wireless
internet-enabled device\\
\rlab{Mobile advertising} & The paid promotional subset of mobile marketing\\
\rowcolor{tablealt}
\rlab{Mobile marketing} & A multi-channel strategy reaching users on mobile
devices\\
\rlab{Push notification} & A short message sent by an app to the device screen,
even when the app is closed\\
\rowcolor{tablealt}
\rlab{Responsive design} & One layout that adapts fluidly to mobile, tablet and
desktop\\
\rlab{Satisficing} & Choosing the first option that plausibly works rather than
the best one\\
\rowcolor{tablealt}
\rlab{UI / UX} & The visual interactive surface / the whole quality of the user's
journey\\
\rlab{User-centred design} & Design driven by explicit user understanding and
validated iteratively by user testing\\
\rowcolor{tablealt}
\rlab{Visual hierarchy} & Making important elements visually prominent in
proportion to their importance\\
\rlab{Web design} & Creating visually appealing, user-friendly websites --- layout,
colour, typography, graphics, navigation, responsiveness\\
\rowcolor{tablealt}
\rlab{Web hosting} & Renting the server that stores the site's files\\
\rlab{Website copy} & The written content of a page --- headlines, body, microcopy,
CTA text\\
\rowcolor{tablealt}
\rlab{Website optimization} & Improving website performance, speed, usability and
search visibility to increase traffic and conversions\\
\end{mmlongtable}

\chapter{Framework Quick-Reference Sheet}

Everything in this volume that is a numbered list, in one place, for the night
before the examination.

\begin{mmlongtable}{Framework quick reference}{tab:quickref4}%
{@{}P{0.28\textwidth}P{0.65\textwidth}@{}}
\rowcolor{ocre}\thd{Framework} & \thd{The sequence}\\
\midrule
\endfirsthead
\rowcolor{ocre}\thd{Framework} & \thd{The sequence}\\
\midrule
\endhead
\rlab{Web design defined} & Layout, colour, typography, graphics, navigation,
responsiveness\\
\rowcolor{tablealt}
\rlab{UI against UX} & Definition, scope, deliverables, marketing role, failure
symptom\\
\rlab{Positive user journey (eight actions)} & Map the journey, shallow
navigation, obvious primary action, fewest steps, feedback, error and empty
states, speed, test with real users\\
\rowcolor{tablealt}
\rlab{Basic pages} & Home, about, product and service, landing, blog, contact,
trust\\
\rlab{Website optimization --- importance} & Speed, experience, ranking, bounce,
conversion, mobile compatibility\\
\rowcolor{tablealt}
\rlab{Four types of optimization} & Technical, content, SEO, mobile\\
\rlab{Eight optimisation techniques} & Speed \ra{} responsive design \ra{}
navigation \ra{} content quality \ra{} strong CTAs \ra{} SEO \ra{} bounce
reduction \ra{} metric monitoring\\
\rowcolor{tablealt}
\rlab{Image optimization --- four gains} & Speed, mobile experience, SEO,
accessibility\\
\rlab{Testing methods} & A/B, usability, heatmaps, recordings, funnel analysis,
performance, cross-device, accessibility, crawl\\
\rowcolor{tablealt}
\rlab{Nine publishing steps} & Requirement analysis \ra{} design \ra{} content
\ra{} domain \ra{} hosting \ra{} development \ra{} testing \ra{} publishing \ra{}
maintenance\\
\rlab{UCD --- seven principles} & Explicit user understanding, continuous
involvement, evaluation-driven, iterative, whole experience, multidisciplinary,
accessible\\
\rowcolor{tablealt}
\rlab{UCD process} & Understand context \ra{} specify requirements \ra{} produce
solutions \ra{} evaluate \ra{} iterate\\
\rlab{Ten visitor questions} & Right place, what's in it for me, trust, cost, next
step, findability, speed, mobile, human contact, safety\\
\rowcolor{tablealt}
\rlab{Twelve design principles} & Simplicity, consistency, visual hierarchy,
balance and alignment, contrast, whitespace, typography, colour, responsiveness,
accessibility, navigation clarity, speed\\
\rlab{Seven usability attributes} & Useful, learnable, memorable, effective,
efficient, desirable, delightful\\
\rowcolor{tablealt}
\rlab{Billboard design rules} & Visual hierarchy, omit needless words, obvious
clickables, breadcrumbs and street signs\\
\rlab{Copy principles} & Benefit-led, scannable, front-loaded, customer's
language, specific, active, concise, one idea per section, explicit CTA, objections
answered in place\\
\rowcolor{tablealt}
\rlab{Design serving copy} & Hierarchy, typography, whitespace, contrast,
consistency, imagery, responsiveness\\
\rlab{Four metrics to name} & Traffic, bounce rate, conversion rate, average
session duration\\
\rowcolor{tablealt}
\rlab{Data to decision} & Observe \ra{} segment \ra{} diagnose \ra{} hypothesise
\ra{} test \ra{} decide \ra{} re-measure\\
\rlab{Mobile advertising against marketing} & Scope, nature, channels, objective,
duration, interaction, measurement --- advertising is the paid subset\\
\rowcolor{tablealt}
\rlab{Mobile web against app} & Access, reach, discoverability, engagement depth,
push, offline, device features, cost\\
\rlab{Five mobile limitations} & Platform diversity, privacy, small-screen
navigation, ad fatigue, regulation\\
\rowcolor{tablealt}
\rlab{Seven campaign characteristics} & Mobile first, multi-screen, full spectrum,
integrate with traditional, work across screens, whole funnel, test your way\\
\rlab{Push uses} & Immediacy, urgency, personalisation, cart recovery,
geo-triggered, re-engagement, loyalty\\
\rowcolor{tablealt}
\rlab{Eight push strategies} & Segment, time correctly, cap frequency, one benefit
with a verb, rich media, deep-link, ask after value, A/B test and watch opt-out\\
\rlab{Ten promotion techniques} & SMS, coupons, in-app exclusives, push flash
sales, location offers, QR codes, gamification, wallet passes, referrals, early
access\\
\rowcolor{tablealt}
\rlab{Mobile ad formats} & Click-to-call, interstitial, click-to-map, canvas and
expandable, mobile search, in-app banner and native, rewarded video\\
\rlab{Eight reasons apps work} & Permanent presence, push access, behavioural
data, higher engagement, loyalty, device capability, offline, frictionless
repurchase\\
\rowcolor{tablealt}
\rlab{Three app rules} & Solve a problem, get inside the user's mind, design with
the end in mind\\
\rlab{Seven mobile metric groups} & Acquisition, activation, engagement,
retention, monetisation, notification performance, technical\\
\end{mmlongtable}
